\section{Phasor approximation and reference frames}\label{sec-refframe}

Under the phasor approximation, the network equations can be written in compact form as:
\begin{equation}\label{iyv}
\bar{\bf I} = {\bf Y} \bar{\bf V}
\end{equation}
where:
\begin{itemize}
\item[$\bar{\bf I}$] is the vector of complex currents injected into the network at the various buses
\item[$\bar{\bf V}$] is the vector of complex voltages at the various buses
\item[$\bf Y$] is the bus (or nodal) admittance matrix.
\end{itemize}

Which frequency consider for the phasors and the admittance matrix ? In dynamic regime, each synchronous machine defines a local frequency. In most cases, those various frequencies remain close to the nominal frequency $f_N$. For large deviations with respect to the nominal value $f_N$, it can be envisaged to update the entries of the $\bf Y$ matrix using the average system frequency, for instance. Otherwise, the admittances are simply computed at frequency $f_N$.

Under the phasor approximation, the voltage at the $i$-th bus takes on the form:
\begin{eqnarray}
v_i(t) & = & \sqrt{2} V_i(t) \cos(\omega_N t + \phi_i(t)) = \sqrt{2} \; re \left[ V_i(t) e^{\,j \, \phi_i(t)} e^{\,j \, \omega_N t}\right] \nonumber\\
       & = & \sqrt{2} \; re \left[ (v_{xi}(t) + j v_{yi}(t)) \; e^{\,j \, \omega_N t} \right] \label{vioft}
\end{eqnarray}
where $v_{xi}(t) + j \, v_{yi}(t)$ is the voltage phasor, in rectangular coordinates, expressed
with respect to $(x,y)$ axes rotating at the angular speed $\omega_N = 2 \pi f_N$, as illustrated in Fig.~\ref{fig-ref_and_init-phasor}.

\begin{figure}[htb]
\centering
\includegraphics[scale=1.0]{phasor.pdf}
\caption{Reference axes and voltage phasor}\label{fig-ref_and_init-phasor}
\end{figure}

Equations similar to \eqref{vioft} can be written for the currents injected into the network, the currents flowing into the network branches, etc.

Instead of determining the ``full wave" evolution of voltages or currents, dynamic simulation in phasor mode aims at rendering the time evolution of $v_{xi}$ and $v_{yi}$\footnote{In RAMSES the rectangular components have been preferred to the polar components $V_i$ and $\phi_i$}.

The $(x,y)$ axes rotating at angular speed $\omega_N$ make up a {\em synchronous reference}.

Although simple, this reference frame suffers from a significant limitation. Indeed, assuming that the power system initially operates at the nominal franquency $f_N$\footnote{This is a common assumption in power system simulation software}, after a disturbance, it will settle at a different frequency $f$, unless its model includes an infinite bus imposing the frequency $f_N$. From there on, the various voltage and current phasors rotate at the angular speed $2 \pi f \ne \omega_N$. Hence, phasor components such as $v_{xi}$ and $v_{yi}$ will oscillate at a frequency $|f-f_N|$, although the system is at equilibrium from a practical viewpoint. For that reason, the synchronous reference is not suitable for long-term simulations, since tracking the oscillations at frequency $|f-f_N|$ requires using a small enough time step size. The synchronous reference frame is suited to short-term simulations (where frequency has not yet returned to steady state) or when the model includes an infinite bus driving the frequency back to $f_N$.

In fact, any (more convenient) speed can be considered for the reference axes $(x,y)$. The only constraint is that all voltage and current phasors refer to the same axes. 

In the {\em Center Of Inertia (COI)} reference frame, the $(x, y)$ axes rotate at the angular frequency:
\begin{equation}\label{omegacoi}
\omega_{coi} = \frac{\sum_{i=1}^m M_i \omega_i}{\sum_{i=1}^m M_i}
\end{equation}
where:
\begin{itemize}
\item[$m$] is the total number of synchronous machines
\item[$\omega_i$] is the rotor speed of $i$-th synchronous machine ($i=1,\ldots,m$)
\item[$M_i$] is the inertia coefficient of the $i$-th machine ($i=1,\ldots,m$).
\end{itemize}
The $M_i$ values relate to the inertia constants $H_i$ (in s) of the individual machines, expressed on a common base power $S_B$ (in MVA):
\[ M_i = 2 H_i \frac{S_{Ni}}{S_B} \]
where $S_{Ni}$ is the nominal apparent power of the $i$-th machine (in MVA).

The COI reference frame does not suffer from the above mentioned drawback. Indeed, when the system settles at a frequency $f$, all synchronous machines rotate at the angular speed $2 \pi f$, and so do the reference axes ($\omega_{coi} = 2 \pi f$). Hence, phasor components such as $v_{xi}$ and $v_{yi}$ tend to constant values, a larger time step size can be used and the simulation is computationally less demanding. The COI reference frame is well suited long-term simulations.

The COI frequency can be used as an average system frequency in injector and two-port models: see Section~\ref{sec-states}.


\subsection{Specifying the reference frame}

The reference to be used is specified in the simulation settings: see Chapter \ref{ch-settings}. 

The presence of a Th\'evenin equivalent (involving a constant frequency voltage source) among the components causes the simulation to use the synchronous reference frame.

\subsection{Implementation of COI reference frame}

To preserve the ``sparsity'' of the model in spite of Eq.~(\ref{omegacoi}), which embeds the rotor speeds of all synchronous machines, the value of $\omega_{coi}$ at the previous time step is used. See:
\begin{quote}
\rm D. Fabozzi and T. Van Cutsem ``On angle references in long-term time-domain simulations'', IEEE Transactions on Power Systems, Vol. 26, No 1, pp. 483-484, Feb. 2011
\end{quote}

for a more detailed presentation of this technique.


%------------------------------------------------------------------------------------------
\section{Network equations}

With all voltage and current phasors referred to the $(x,y)$ axes, the network equations can be decomposed into:
\begin{equation}
{\bf i}_x + j \, {\bf i}_y = {\bf Y} \left( {\bf v}_x + j \, {\bf v}_y \right) = 
\left( {\bf G} + j \, {\bf B} \right) \left( {\bf v}_x + j \, {\bf v}_y \right)
\end{equation}
with:
\begin{equation}
{\bf v}_x = \left[ \begin{array}{c} v_{x1} \\ \vdots \\ v_{xN}
\end{array} \right] \;\;\;\;\;\; {\bf v}_y = \left[ \begin{array}{c} v_{y1} \\ \vdots \\ v_{yN} \end{array} \right] \;\;\;\;\;\; {\bf i}_x = \left[ \begin{array}{c} i_{x1} \\ \vdots \\ i_{xN} \end{array} \right] \;\;\;\;\;\; {\bf i}_y = \left[ \begin{array}{c} i_{y1} \\ \vdots \\ i_{yN}
\end{array}\right]
\end{equation}
where ${\bf G}$ is the conductance matrix and ${\bf B}$ the susceptance matrix.

Decomposing into real and imaginary parts and assembling into a single equation yields:
\begin{equation}\label{neteqs}
\left[ \begin{array}{c} {\bf i}_x \\ {\bf i}_y \end{array} \right] = \left[ \begin{array}{cc} {\bf G} & - {\bf B} \\ {\bf B} & {\bf G} \end{array} \right] \left[ \begin{array}{c} {\bf v}_x \\ {\bf v}_y \end{array} \right]
\end{equation}
Thus, for a network with $N$ buses, there are $2N$ equations \eqref{neteqs} involving $4N$ variables.

%------------------------------------------------------------------------------------------
\section{Initialization procedure}\label{sec-ini-proc}

The dynamic simulation is initialized according to the procedure detailed next.
\begin{enumerate}
\item The starting point is the vector of initial bus voltages, as detailed in Section~\ref{sec-lfresv}
\item setting the bus voltages to those values, the active and reactive power flows in the network branches and shunts are computed
\item by summing the power flows in incident branches, the bus power injections are determined. This is illustrated in Fig.~\ref{fig-bus-balance}, where $P^{inj}_i$ (resp. $Q^{inj}_i$) is the active (resp. reactive) power injected into the $i$-th bus. A \emph{positive value} corresponds to power \emph{entering the network}
\item at each bus, the bus power injection is shared among the various components (generators, loads, injectors, etc.) connected to that bus. 

There are two ways of assigning power to the $j$-th component:
\begin{itemize}
\item[(i)] by specifying the active and reactive powers injected into the network by that component:
\begin{equation}\label{inipower}
P^{c}_j = P^{c0}_j \:\:\:\:\:\:\:\:\:\: Q^{c}_j = Q^{c0}_j
\end{equation}
\item[(ii)] by specifying the fraction of the bus power injection taken by that component :
\begin{equation}\label{frac}
P^{c}_j = f_{Pj} \; P^{inj}_i \:\:\:\:\:\:\:\:\:\: Q^{c}_j = f_{Qj} \; Q^{inj}_i
\end{equation}
where $f_{Pj}$ (resp. $f_{Qj}$) denotes the fraction relative to active (resp. reactive) power.
\end{itemize}

\item the remaining power, not taken by the components at bus $i$, i.e.
\[ P_i^r = P^{inj}_i - \sum_{j=1}^n P^c_j \:\:\:\:\:\:\:\:\:\: Q_i^r = Q^{inj}_i - \sum_{j=1}^n Q^c_j \]
is checked. If larger than an internal tolerance, the power is assigned to an impedance load\footnote{to be more precise a constant shunt impedance, or equivalently a constant shunt admittance load}, as shown in Fig.~\ref{fig-bus-balance}. The load admittance is such that the power balance is satisfied at $t=0$:
\begin{equation}\label{ini-adm}
(G_i^r - j \: B_i^r) \; V_i(0)^2 = - (P_i^r + \: j Q_i^r)
\end{equation}
where $V_i(0)$ is the initial bus voltage magnitude\footnote{the left-hand side of Eq.~(\ref{ini-adm}) is the complex power consumed by the admittance, while $P_i^r + \: j Q_i^r$ is counted positive when the power is injected into the network. Hence, the minus sign in the right hand side of the equation}. Note that the sign of $G_i^r$ is opposite to that of $P_i^r$, while the sign of $B_i^r$ is the same as that of $Q_i^r$.
\end{enumerate}

\begin{figure}[htb]
\centering
\includegraphics[scale=1.0]{init-pflow.pdf}
\caption{initial powers at bus $i$}\label{fig-bus-balance}
\end{figure}

The two ways  mentioned under items (i) and (ii) above are \emph{mutually exclusive}, i.e. either a nonzero power or a nonzero fraction is specified, but not both. In mathematical terms:
\begin{equation}\label{constr-ini-data}
f_{Pj} \times P^{c0}_j = 0  \:\:\:\:\:\:\:\:\:\: f_{Qj} \times Q^{c0}_j = 0
\end{equation}

Although they are typically in the interval $[0 \; 1]$, the values of the $f_{Pj}$'s and $f_{Qj}$'s may be negative or larger than one.

In the quite common case when a single component is connected to a bus, and it takes the whole power consumed or produced at that bus, it is convenient to specify:
\[ f_{Pj} = 1 \:\:\:\:\:\:\:\:\:\:\:\:  P^{c0}_j = 0  \:\:\:\:\:\:\:\:\:\:\: f_{Qj} = 1 \:\:\:\:\:\:\:\:\:\:\: Q^{c0}_j = 0 \]
With fractions instead of powers, the data can be re-used at another initial operating point\footnote{i.e. when another set of initial voltages is used}.

The constant admittance (or impedance) loads corresponding to Eq.~(\ref{ini-adm}) are created automatically at system initialization. They are given names of the type {\tt M\_bus}, where {\tt M} stands for ``mismatch'' and {\tt bus} is the name of the bus to which the admittance is connected. Names starting with {\tt M\_} are reserved and must not be used when defining constant admittance loads.

A large value of $G_i^r$ (resp. $B_i^r$) may be intentional, when a load is to be modelled as a constant shunt admittance and the task is left to RAMSES. However, it may also result from a mistake in the initial power balance. It is recommended to check the values of the {\tt M\_} loads through the menu ``Load Initialization'' in the STEPSS interface.

\vspace*{-2ex}
\begin{quote}
\faHandPointRight ~ Here is an example of output produced by RAMSES at initialization.
\vspace*{-2ex}
\begin{verbatim}
NUMBER OF IMPEDANCE LOADS :     3  (M_ type:     3 )

load name              bus name                  P          Q

M_2                    2                       90.002     17.997
M_3                    3                        0.013     -0.011
M_4                    4                       -0.017     -0.022

\end{verbatim}
In this example, three {\tt M\_} loads have been produced at initialization, respectively at buses 2, 3 and 4. The one at bus 3 consumes 0.013 MW and produces 0.011 Mvar. This load was created because it is larger than the internal tolerance but, for a transmission system, this is a negligible value. The same holds true for the {\tt M\_} load at bus 4. On the other hand, the {\tt M\_} load at bus 2 is non negligible. This could be due to forgetting to specify a load at bus 2.
\end{quote}

%------------------------------------------------------------------------------------------
\section{Data format}

The way to specify the values of $f_{Pj}, P^{c0}_j, f_{Qj}$ and $Q^{c0}_j$ is detailed in the following subsequent sections of this documentation:
\begin{itemize}
\item for a synchronous machine: Section~\ref{format-sync}
\item for a Th\'{e}venin equivalent: Section~\ref{format-thev}
\item for a constant impedance load: Section~\ref{format-impload}
\item for an injector with user-defined model: Section~\ref{format-inj}
\item for a two-port with user-defined model: Section~\ref{format-twop}
\end{itemize}


 